% GENERATED FILE. Edit canonical lessons, then run npm run generate:books.
% canonical-source-hash: fnv1a64:97e53cc2b191584e
% canonical-lessons: PT-C487-sintoma, PT-C487-analises, PT-C487-diagnostico, PT-C487-receitar, PT-C487-tratamento

\chapter{Na consulta}
\label{ch:pt-na-consulta}
\hlchaptermodality{487}

\begin{chapteropening}
\textbf{By the end of this chapter:} \emph{I can describe symptoms, tests and treatment at the doctor.}

Complete the last lesson of chapter 487: I can describe symptoms, tests and treatment at the doctor.
\end{chapteropening}

\section[o sintoma]{o sintoma — symptom}

\label{lesson:PT-C487-sintoma}



\begin{quote}
Before the new one: say the Portuguese for sprain, then the Portuguese for bandage.
\end{quote}

\subsection*{You'll want to know: o sintoma}
\textbf{o sintoma} — \textquotedblleft{}symptom\textquotedblright{}.

\begin{quote}
A febre e a tosse são sintomas da gripe.
\end{quote}

\begin{quote}
\emph{Fever and a cough are symptoms of flu.}
\end{quote}

Masculine, though it ends in \emph{-a}: \emph{o sintoma}.

\begin{cousinweb}
\emph{sintoma} goes back through Latin to Greek \emph{sýmptōma} \textquotedblleft{}a happening, a chance event, a symptom\textquotedblright{} (EtymDB).
\end{cousinweb}

\begin{grammarlens}[title={What you've built}]
One more word for telling, explaining and describing at B1.
\end{grammarlens}

\subsection*{Your turn}
\begin{itemize}
\raggedright
  \item \emph{Say it:} \emph{o sintoma}
  \item \emph{Say it:} \emph{o sintoma}, once more
  \item \emph{Say it:} a sentence of your own with \emph{o sintoma}
  \item \emph{Recall:} say the Portuguese for sprain, then the Portuguese for bandage, then say \emph{o sintoma} again
\end{itemize}

\begin{tcolorbox}[breakable,colback=teal!4,colframe=teal!35!black,title={Before you move on}]
What does o sintoma mean? (\textquotedblleft{}Symptom\textquotedblright{}.) Say it once more.
\end{tcolorbox}
\section[as análises]{as análises — (medical) tests}

\label{lesson:PT-C487-analises}



\begin{quote}
Before the new one: say the Portuguese for bandage, then the Portuguese for symptom.
\end{quote}

\subsection*{You'll want to know: as análises}
\textbf{as análises} — \textquotedblleft{}(medical) tests\textquotedblright{}.

\begin{quote}
O médico mandou-me fazer análises ao sangue.
\end{quote}

\begin{quote}
\emph{The doctor sent me for blood tests.}
\end{quote}

\emph{Fazer análises} — to have tests done. Singular \emph{a análise} — an analysis.

\begin{cousinweb}
\emph{análise} is a learned word, usually traced through Medieval Latin to Greek \emph{análysis} \textquotedblleft{}a loosening, a breaking up\textquotedblright{}.
\end{cousinweb}

\begin{grammarlens}[title={What you've built}]
One more word for telling, explaining and describing at B1.
\end{grammarlens}

\subsection*{Your turn}
\begin{itemize}
\raggedright
  \item \emph{Say it:} \emph{as análises}
  \item \emph{Say it:} \emph{as análises}, once more
  \item \emph{Say it:} a sentence of your own with \emph{as análises}
  \item \emph{Recall:} say the Portuguese for bandage, then the Portuguese for symptom, then say \emph{as análises} again
\end{itemize}

\begin{tcolorbox}[breakable,colback=teal!4,colframe=teal!35!black,title={Before you move on}]
What does as análises mean? (\textquotedblleft{}(medical) tests\textquotedblright{}.) Say it once more.
\end{tcolorbox}
\section[o diagnóstico]{o diagnóstico — diagnosis}

\label{lesson:PT-C487-diagnostico}



\begin{quote}
Before the new one: say the Portuguese for symptom, then the Portuguese for (medical) tests.
\end{quote}

\subsection*{You'll want to know: o diagnóstico}
\textbf{o diagnóstico} — \textquotedblleft{}diagnosis\textquotedblright{}.

\begin{quote}
O diagnóstico foi uma simples constipação.
\end{quote}

\begin{quote}
\emph{The diagnosis was a simple cold.}
\end{quote}

\emph{Fazer o diagnóstico} — to make the diagnosis.

\begin{cousinweb}
\emph{diagnóstico} is a learned word, usually traced to Greek \emph{diagnōstikós} \textquotedblleft{}able to tell apart\textquotedblright{}, built on \emph{gign\'{\={o}}skein} \textquotedblleft{}to know\textquotedblright{}.
\end{cousinweb}

\begin{grammarlens}[title={What you've built}]
One more word for telling, explaining and describing at B1.
\end{grammarlens}

\subsection*{Your turn}
\begin{itemize}
\raggedright
  \item \emph{Say it:} \emph{o diagnóstico}
  \item \emph{Say it:} \emph{o diagnóstico}, once more
  \item \emph{Say it:} a sentence of your own with \emph{o diagnóstico}
  \item \emph{Recall:} say the Portuguese for symptom, then the Portuguese for (medical) tests, then say \emph{o diagnóstico} again
\end{itemize}

\begin{tcolorbox}[breakable,colback=teal!4,colframe=teal!35!black,title={Before you move on}]
What does o diagnóstico mean? (\textquotedblleft{}Diagnosis\textquotedblright{}.) Say it once more.
\end{tcolorbox}
\section[receitar]{receitar — to prescribe}

\label{lesson:PT-C487-receitar}



\begin{quote}
Before the new one: say the Portuguese for (medical) tests, then the Portuguese for diagnosis.
\end{quote}

\subsection*{You'll want to know: receitar}
\textbf{receitar} — \textquotedblleft{}to prescribe\textquotedblright{}.

\begin{quote}
O médico receitou-me um xarope.
\end{quote}

\begin{quote}
\emph{The doctor prescribed me a syrup.}
\end{quote}

Built on \emph{a receita}, which is a prescription and also a recipe.

\begin{cousinweb}
\emph{receitar} is built on \emph{receita} \textquotedblleft{}a prescription, a recipe\textquotedblright{}, from Latin \emph{recepta} \textquotedblleft{}things received\textquotedblright{} (EtymDB), from \emph{recipere} \textquotedblleft{}to take back\textquotedblright{} (Lewis \& Short).
\end{cousinweb}

\begin{grammarlens}[title={What you've built}]
One more word for telling, explaining and describing at B1.
\end{grammarlens}

\subsection*{Your turn}
\begin{itemize}
\raggedright
  \item \emph{Say it:} \emph{receitar}
  \item \emph{Say it:} \emph{receitar}, once more
  \item \emph{Say it:} a sentence of your own with \emph{receitar}
  \item \emph{Recall:} say the Portuguese for (medical) tests, then the Portuguese for diagnosis, then say \emph{receitar} again
\end{itemize}

\begin{tcolorbox}[breakable,colback=teal!4,colframe=teal!35!black,title={Before you move on}]
What does receitar mean? (\textquotedblleft{}To prescribe\textquotedblright{}.) Say it once more.
\end{tcolorbox}
\section[o tratamento]{o tratamento — treatment}

\label{lesson:PT-C487-tratamento}



\begin{quote}
Before the new one: say the Portuguese for diagnosis, then the Portuguese for to prescribe.
\end{quote}

\subsection*{You'll want to know: o tratamento}
\textbf{o tratamento} — \textquotedblleft{}treatment\textquotedblright{}.

\begin{quote}
O tratamento dura duas semanas.
\end{quote}

\begin{quote}
\emph{The treatment lasts two weeks.}
\end{quote}

Built on \emph{tratar} — to treat.

\begin{cousinweb}
\emph{tratamento} is built on \emph{tratar} \textquotedblleft{}to treat\textquotedblright{}, from Latin \emph{tractāre} \textquotedblleft{}to handle\textquotedblright{} (Lewis \& Short), from \emph{trahere} \textquotedblleft{}to draw\textquotedblright{}.
\end{cousinweb}

\begin{grammarlens}[title={What you've built}]
One more word for telling, explaining and describing at B1.
\end{grammarlens}

\subsection*{Your turn}
\begin{itemize}
\raggedright
  \item \emph{Say it:} \emph{o tratamento}
  \item \emph{Say it:} \emph{o tratamento}, once more
  \item \emph{Say it:} a sentence of your own with \emph{o tratamento}
  \item \emph{Recall:} say the Portuguese for diagnosis, then the Portuguese for to prescribe, then say \emph{o tratamento} again
\end{itemize}

\begin{tcolorbox}[breakable,colback=teal!4,colframe=teal!35!black,title={Before you move on}]
What does o tratamento mean? (\textquotedblleft{}Treatment\textquotedblright{}.) Say it once more.
\end{tcolorbox}
